\documentclass[11pt]{article}

\usepackage[a4paper,margin=1in]{geometry}
\usepackage[T1]{fontenc}
\usepackage{lmodern}
\usepackage{microtype}
\usepackage{amsmath,amssymb,amsthm,mathtools}
\usepackage{booktabs,tabularx}
\usepackage{enumitem}
\usepackage[numbers,sort&compress]{natbib}
\usepackage[colorlinks=true,linkcolor=blue,citecolor=blue,urlcolor=blue]{hyperref}
\usepackage[nameinlink,noabbrev]{cleveref}

\setlength{\emergencystretch}{3em}

\newtheorem{theorem}{Theorem}[section]
\newtheorem{proposition}[theorem]{Proposition}
\newtheorem{corollary}[theorem]{Corollary}
\newtheorem{lemma}[theorem]{Lemma}
\theoremstyle{definition}
\newtheorem{definition}[theorem]{Definition}
\theoremstyle{remark}
\newtheorem{remark}[theorem]{Remark}

\newcommand{\Lzero}{\mathrm{L0}}
\newcommand{\Lone}{\mathrm{L1}}
\newcommand{\Ltwo}{\mathrm{L2}}

\title{\textbf{Quantitative Shift-One Shadows of the
Determinant-Two Frame}\\
\large Exact Even-Carrier Transfer and the Active Odd-Carrier
Firewall}
\author{Liang Wang}
\date{July 2026}

\begin{document}
\maketitle

\begin{abstract}
The unrestricted shift-two Liouville sign has an exact quantitative
shadow on its even carrier:
\[
 \lambda(2m-2)\lambda(2m)=\lambda(m-1)\lambda(m).
\]
Consequently, proved quantitative logarithmic and almost-all-scale
results for the consecutive shift-one correlation transfer to that
even subcarrier.  This does not estimate the active
determinant-two TPC packet.  On the simultaneous squarefree support
of TPC-127, the affine values \(D,V\) are odd, and hence
\(n=sV=aD+2\) is odd.  We prove a complementary determinant
invariant: on any residue progression contained in the odd
integers, extracting all fixed contents from \(n-2,n\) leaves an
affine determinant divisible by four.  It cannot produce the
determinant-one pair \(t,t+1\).  We cite
Helfgott--Radziwi{\l}{\l} and Pilatte only for their stated
shift-one conclusions.  Separately, Tao--Ter\"av\"ainen now provide
a quantitative almost-all-scale theorem uniformly for positive
affine data of sufficiently small polylogarithmic height.  That
new corridor can cover eligible determinant-two affine components,
but it does not supply the squarefree component reassembly,
deterministic all-prefix selector, or physical return.  The transfer
and firewall are exact \(\Lzero\) statements; the full actual packet
remains open, and no positive \(\Ltwo\) result is claimed.
\end{abstract}

\section{Four theorem scopes that must not be merged}

Let \(\lambda(n)=(-1)^{\Omega(n)}\).  There are three distinct
claims in play.

\paragraph{Fixed affine, qualitative logarithmic.}
Tao proved logarithmically averaged two-point cancellation for
arbitrary fixed nonparallel affine forms \citep{Tao2016LogChowla}.
TPC-137 uses this theorem to close the full squarefree masks for
fixed determinant-two data \citep{WangTPC137}.

\paragraph{Consecutive shift one, quantitative.}
Helfgott and Radziwi{\l}{\l} proved, among other statements, a
quantitative logarithmic bound for
\(\lambda(n)\lambda(n+1)\) and ordinary cancellation at almost all
scales \citep{HelfgottRadziwill2021}.  Pilatte later proved
\begin{equation}
 \sum_{n\le x}\frac{\lambda(n)\lambda(n+1)}n
 \ll(\log x)^{1-c}
\label{eq:pilatte}
\end{equation}
for some absolute \(c>0\) \citep{Pilatte2023}.

\paragraph{Small-polylog affine, quantitative almost all scales.}
Tao and Ter\"av\"ainen prove the following uniform special case of
their quantitative correlation theorem: for some sufficiently
small absolute \(c_0>0\), outside a set of scales of logarithmic
density \(O((\log X)^{-c_0})\),
\begin{equation}
 \frac1N\sum_{n\le N}
 \lambda(a_1n+b_1)\lambda(a_2n+b_2)
 \ll(\log N)^{-c_0}
\label{eq:tt-affine}
\end{equation}
uniformly for positive integers
\[
 a_1,a_2,b_1,b_2\le(\log N)^{c_0},
 \qquad a_1b_2\ne a_2b_1
\]
\citep[Theorem~3.1 and Remarks~3.2]{TaoTeravainen2026}.
Shrinking \(c_0\) absorbs harmless differences between the height
and saving exponents.  This is a genuine arbitrary-affine
quantitative corridor, but only at nonexceptional scales and only
inside its stated small-polylog range.

\paragraph{Actual determinant two.}
TPC-127 produces
\[
 n=sV(z)=aD(z)+2
\]
on one generally growing progression, behind quotient-squarefree
masks, physical weights, phases and prefixes
\citep{WangTPC127}.  Neither \eqref{eq:pilatte} nor the stated
shift-one almost-scale result is, as cited here, a theorem on this
actual family.  Equation \eqref{eq:tt-affine} can estimate an
individual reduced component only after its literal coefficients,
constants and CRT modulus have been proved to lie in the
small-polylog corridor; it still leaves exceptional-scale
selection and complete reassembly.

\begin{definition}[Source-safe quantitative gate]
\label{def:gate}
A quantitative determinant-two affine theorem is a separately
proved estimate whose statement itself includes the desired
determinant-two affine forms and all required growth quantifiers.
A shift-one theorem plus an unproved assertion that its method
``should extend'' does not pass this gate.  The proved
Tao--Ter\"av\"ainen corridor passes only for records satisfying its
explicit height and nonexceptional-scale hypotheses.
\end{definition}

\section{The exact even-carrier transfer}

\begin{proposition}[Shift two becomes shift one on even integers]
\label{prop:even}
For every \(m\ge2\),
\begin{equation}
\boxed{
 \lambda(2m-2)\lambda(2m)
 =
 \lambda(m-1)\lambda(m).}
\label{eq:even}
\end{equation}
Consequently, for any finite interval \(J\),
\begin{equation}
 \sum_{\substack{n\in2J\\n\ {\rm even}}}
 \frac{\lambda(n-2)\lambda(n)}n
 =
 \frac12\sum_{m\in J}
 \frac{\lambda(m-1)\lambda(m)}m.
\label{eq:harmonic-transfer}
\end{equation}
The corresponding unweighted interval sums agree after the same
change of variable, up to endpoint conventions.
\end{proposition}

\begin{proof}
Complete multiplicativity gives
\[
 \lambda(2m-2)\lambda(2m)
 =
 \lambda(2)^2\lambda(m-1)\lambda(m),
\]
and \(\lambda(2)^2=1\).  Substitute \(n=2m\) to obtain the two sum
identities.
\end{proof}

\begin{corollary}[What the quantitative source theorems control]
\label{cor:shadow}
Every stated shift-one logarithmic or almost-all-scale estimate
transfers, with its native quantifiers and rate, to the even
subcarrier of the unrestricted sequence
\(\lambda(n-2)\lambda(n)\).
\end{corollary}

\begin{proof}
Apply \cref{prop:even}; the fixed dilation by two changes only
absolute constants and interval endpoints.
\end{proof}

This is genuine arithmetic information, but its carrier must be
checked before it is inserted in H3.

\section{The actual squarefree carrier is odd}

TPC-127 proves the following arithmetic geometry.  If
\[
 D(z)=d+sz,\qquad V(z)=u+az,\qquad su-ad=2,
\]
with \(a,s\) odd, then \(D(z),V(z)\) have the same parity.  On their
simultaneous squarefree support, the even--even case is impossible.

\begin{proposition}[Empty overlap]
\label{prop:empty}
On the simultaneous squarefree determinant-two support,
\[
 D(z)\equiv V(z)\equiv n(z)\equiv1\pmod2,
\qquad n(z)=sV(z)=aD(z)+2.
\]
Hence the carrier controlled by \cref{cor:shadow} has empty
intersection with the active squarefree carrier.
\end{proposition}

\begin{proof}
The oddness and coprimality of \(D,V\) are the parity classification
of TPC-127.  Since \(s\) is odd, \(n=sV\) is odd.  The carrier in
\cref{cor:shadow} consists of even \(n\).
\end{proof}

\begin{remark}[The positive result is not discarded]
\Cref{prop:empty} does not weaken the shift-one theorems.  It says
that their exact even-carrier consequence bounds a different
subsequence.  Calling it an estimate of the odd squarefree carrier
would be a support error.
\end{remark}

\section{An odd-carrier determinant invariant}

The possibility remains that a residue change and extraction of
fixed contents might convert the odd carrier to consecutive forms.
The next theorem rules out that elementary move.

\begin{theorem}[Even determinant survives odd-carrier reduction]
\label{thm:invariant}
Let \(M\) be even and \(r\) odd, so that
\(n=r+Mt\) is odd for every integer \(t\).  Put
\[
 c_-=(M,r-2),\qquad c_+=(M,r),
\]
and define the primitive affine forms
\[
 L_-(t)=\frac{Mt+r-2}{c_-},
 \qquad
 L_+(t)=\frac{Mt+r}{c_+}.
\]
Then \(c_-,c_+\) are odd and coprime, and
\begin{equation}
\boxed{
 \det(L_-,L_+)=\frac{2M}{c_-c_+}\in4\mathbb Z.}
\label{eq:reduced-det}
\end{equation}
In particular, this pair cannot equal, up to order and fixed
Liouville factors, the primitive consecutive pair \(t+b,t+b+1\),
whose determinant has absolute value \(1\).
\end{theorem}

\begin{proof}
Both \(r\) and \(r-2\) are odd, so both contents are odd.
Furthermore
\[
 (c_-,c_+)\mid(r-2,r)\mid2,
\]
and hence the contents are coprime.  Since each divides \(M\), their
product divides \(M\).  Each displayed affine form is primitive
because its content has been divided out.  The determinant of the
coefficient pairs is
\[
 \frac{M}{c_-}\frac{r}{c_+}
 -
 \frac{M}{c_+}\frac{r-2}{c_-}
 =
 \frac{2M}{c_-c_+}.
\]
Because \(M\) is even and \(c_-c_+\) is odd,
\(M/(c_-c_+)\) is even; hence the determinant is divisible by four.
The determinant of
\((1,b),(1,b+1)\) is \(1\).
\end{proof}

\begin{remark}[Scope of the obstruction]
\Cref{thm:invariant} forbids a residue/content relabeling, not a new
analytic proof.  One may try to rebuild an expansion or entropy
argument directly for determinant two.  That effort is precisely
the open gate in \cref{def:gate}.
\end{remark}

\section{Finite diagnostic and claim boundary}

The deterministic script verifies \eqref{eq:even} term by term,
enumerates simultaneous squarefree determinant-two samples, and
checks primitivity, content coprimality and
\eqref{eq:reduced-det} over many odd residue classes.  These are
exact regressions.  It may also report finite correlations, but no
fitted decay is used as an asymptotic premise.

\begin{center}
\small
\begin{tabularx}{\textwidth}{@{}p{0.48\textwidth}X@{}}
\toprule
Statement & Level and status\\
\midrule
Even shift-two to shift-one transfer
& Exact theorem (\(\Lzero\)).\\
Quantitative consequence on the even carrier
& Proved external input with unchanged shift-one scope.\\
Active determinant-two carrier is odd and disjoint
& Exact support theorem (\(\Lzero\)).\\
Odd-carrier reduced determinant remains even
& Exact firewall (\(\Lzero\)).\\
Quantitative logarithmic or almost-scale theorem on active
determinant two
& Proved for eligible small-polylog affine data outside a
small log-density exceptional scale set; actual-family return
open.\\
Growing physical all-prefix power estimate
& Open \(\Ltwo\) target.\\
\bottomrule
\end{tabularx}
\end{center}

Thus TPC-138 identifies a real but ineligible quantitative shadow.
The next task is not another parity relabeling: it is to compare the
literal squarefree/periodic decomposition with the proved
small-polylog affine corridor, account for all component losses,
and prevent deterministic prefixes from concentrating on its
exceptional scales.

\begingroup
\footnotesize
\raggedright
\bibliographystyle{plainnat}
\bibliography{references}
\endgroup

\end{document}
